% =====================================================================
% Living study report (fork template).
%
% THIS IS A LIVING DOCUMENT. It is updated as the research progresses,
% alongside ProjectProposal.md and docs/. The results section stays empty
% of findings until real data exists; no fabricated results, ever.
%
% Build:  tectonic report.tex        (or: latexmk -pdf report.tex)
% Word:   pandoc report.md -o report.docx   (Markdown body kept in sync)
% =====================================================================
\documentclass[11pt]{article}

\usepackage[margin=1in]{geometry}
\usepackage{amsmath}
\usepackage{graphicx}
\usepackage{booktabs}
\usepackage{hyperref}
\usepackage{enumitem}

\graphicspath{{figures/}}

\title{Living Study Report}
\author{AISoc Researcher}
\date{\today}

\begin{document}
\maketitle

\begin{abstract}
Replace this abstract with a short statement of the problem, the literature gap,
the one research question, the design, and the fact that this is a living report
whose results stay empty until real data is collected. No result is fabricated.
\end{abstract}

\tableofcontents

\section{Introduction}

Replace this section with the background, the problem, and the literature gap
that motivate the study, drawn from \texttt{docs/01-literature-review.md} and
\texttt{docs/02-problem-analysis.md}.

The single main research question of this study is:
\begin{quote}
\emph{<Your one main research question goes here.>}
\end{quote}

State how the question follows from the gap, and what a novel, testable, grounded
answer would look like. Keep this consistent with
\texttt{docs/03-research-question.md} and \texttt{ProjectProposal.md}.

\section{Method}

Replace this section with the study design from
\texttt{docs/05-experiment-design.md} and the pre-registered analysis plan from
\texttt{docs/07-statistical-analysis.md}. State the independent and dependent
variables, the control condition, the outcome metric, the procedure and its
timing, and the statistical model. Describe how the sandboxed runner in
\texttt{experiments/} executes the protocol automatically and records every run
in \texttt{results/ledger.csv}.

\section{Results}

This section stays empty of findings until real data exists. No result is
fabricated. When data is collected through the sandboxed runner or entered by
the researcher, report the outcome metric by condition, the effect, and its
uncertainty here, generated from \texttt{results/} by the scripts in
\texttt{analysis/}, and mirror it into \texttt{docs/06-experimental-results.md}.


\section{Threats to Validity and Conclusion Rules}
The interpretation rules, the four-category threats-to-validity treatment
(internal, construct, external, statistical-conclusion), and the criteria by
which the study reaches a defensible claim are fixed in advance and recorded in
\texttt{docs/08-evaluation.md}. A null result is a valid, reportable outcome
provided the pre-registered analysis in \texttt{docs/07-statistical-analysis.md}
is followed.

\section*{References}
The source index, with reading-depth tags, is maintained in
\texttt{references.md}. The forward reading queue is in \texttt{readingList.md}.

\end{document}
